\section{L2、L3、L4：责任边界比等级数字更重要}

上一章解释能力训练，本章看自动驾驶等级。郎咸朋强调，“L3 不是 L2 的延长，而是 L4 的先导”。这个判断的关键不在数字，而在责任：L2 是辅助驾驶，人类驾驶员负责；L3 在特定条件下系统开始承担责任，并在需要接管时提前通知人；L4 则是在限定区域或限定场景里系统完整负责。

\begin{figure}[H]
\centering
\includegraphics[width=0.92\textwidth]{l3-l4-ladder.png}
\caption{L3 不是 L2 的延长：L3 更像 L4 的先导，而不是辅助驾驶线性升级。自制概念图，依据 00:50:50--01:15:15 对谈内容整理。}
\end{figure}

\begin{knowledgebox}{读图：等级变化背后是责任变化}
左侧 L2 是人负责、系统辅助；右侧 L3/L4 开始要求系统承担责任、明确 ODD 和接管边界。读等级时不要只看功能多不多，而要看事故责任、接管机制、安全冗余和运行边界。
\end{knowledgebox}

\subsection{L2 为什么不能自然长成 L3}

本节解释误区。很多人把 L2、L3、L4 想成连续等级：功能越来越多，体验越来越强。但从责任角度看，L2 到 L3 是质变。L2 可以要求人随时注意并接管；L3 则必须保证在系统负责期间足够安全，并且在退出 ODD 或能力边界时给出可执行的接管窗口。它涉及法规、产品定义、安全策略和用户教育。

\begin{warningbox}{L2 做得越强，用户越容易误用}
如果 L2 体验接近自动驾驶，但法律和安全责任仍要求人负责，用户可能过度信任系统。这是 L2/L3 边界最危险的地方：体验像自动驾驶，责任却还在人。
\end{warningbox}

\subsection{ODD 与接管：责任边界怎样落到产品}

本节把等级问题落到产品设计。L3/L4 不是在发布会上多写一个等级，而是系统必须知道自己能在哪里负责、何时退出、怎样通知人接管、接管失败时怎样进入最小风险状态。例如高速、城市快速路、拥堵跟车、泊车、特定天气和特定速度范围，都可能构成不同 ODD。只要系统承诺在某个 ODD 内负责，它就需要相应的感知冗余、故障诊断、状态监控和法规责任设计。

\begin{importantbox}{责任边界的三个问题}
第一，系统何时可以接管驾驶责任？第二，系统何时必须把责任还给人或进入降级状态？第三，如果人没有及时接管，系统如何安全停车或保持最低风险？这三个问题没有回答清楚，L3 就只是营销词。
\end{importantbox}

\subsection{L3 为什么是 L4 的先导}

本节解释“先导”。L3 要求系统在某些条件下真正负责，这迫使公司建立 L4 需要的能力：ODD 定义、安全冗余、接管策略、风险最小化、系统监控、事故责任和验证体系。因此它不只是 L2 功能包升级，而是进入“系统负责”的训练场。

\noindent\begin{tabularx}{\textwidth}{p{0.16\textwidth}p{0.36\textwidth}X}
\toprule
等级 & 核心责任 & 工程含义 \\
\midrule
L2 & 人负责，系统辅助 & 功能体验可以强，但人必须监督。 \\
L3 & 特定条件下系统负责，需要接管机制 & 责任边界、退出策略和法规认证变关键。 \\
L4 & 限定区域/场景系统完整负责 & 需要完整安全闭环和运行域治理。 \\
L5 & 无限定自动驾驶 & 当前更像远期概念。 \\
\bottomrule
\end{tabularx}

\subsection{本章小结}

自动驾驶等级不是营销数字，而是责任结构。L3 的关键价值，是让公司开始面对 L4 所需的安全、法规、接管和系统责任问题。

